\section{相关工作}
\label{sec:related}

合理反例检验涉及三方面的研究：不确定海流中的航路规划、元推理与计算的价值，以及面向决策的信息获取。

\subsection{不确定海流中的航路规划}
确定性海流中的时间最优航线可由图搜索或水平集方法求得\cite{dijkstra1959note,lolla2014levelset}，区域海洋模式的预报也已用于引导航行器追踪演变中的海洋过程\cite{smith2010planning}。为应对预报的不确定性，研究者发展了集合规划\cite{wang2016path}、基于随机动态正交水平集方程的时间最优规划\cite{subramani2018stochastic}，以及计入预报不确定性的图搜索规划\cite{kularatne2018forecast}；数据同化则用集合卡尔曼滤波等方法把新观测并入海流估计\cite{anderson2001eakf}。这些方法或假设海流图准确，或在给定的集合成员上求期望最优，都没有判断一张海流图何时对某个具体决定不可信。

\subsection{元推理与计算的价值}
元推理研究计算本身的价值，以决定何时停止思考、开始行动\cite{horvitz1987reasoning,russell1991right}；在机器人规划中，这一思想体现为学习停止规则、自适应分配计算预算，以及让规划与执行并行\cite{xiao2020thinking,sung2021quit,lin2015metareasoning,elboher2023formal,magnino2026metareasoning}。然而，在有误差的模型上择优，选中方案的价值往往被高估，即决策分析中的优化者诅咒\cite{smith2006curse}；模型越不准，有效的规划步长越短\cite{jiang2015horizon}。智能体推理的研究也发现，外部状态未知时，增加推理计算不能代替与环境交互\cite{hunt2021formalizing,shen2025thinking,kobalczyk2025active,cuadron2025overthinking}。这些工作决定何时停止计算，却默认继续计算总能改善决策；本文检验的正是这一前提何时不成立。

\subsection{面向决策的信息获取}
只有能区分竞争假设、并可能改变行动的观测才有决策价值\cite{chernoff1959sequential,box1967discrimination}。最优臂识别、信息导向采样与目标导向实验设计都针对具体决策采集样本，而不是均匀降低全局方差\cite{garivier2016optimal,russo2018thompson,russo2018ids,attia2018goal}；面向全场估计的自适应采样则引导航行器前往不确定性高的海区\cite{srinivas2012gpucb,hollinger2014sampling,fossum2019}。鲁棒决策与预测多解性研究进一步主张，当多个与数据相容的模型给出相反结论时，应考察不确定集内的最坏表现\cite{benhaim2006infogap,marx2020predictive}。本文把这一思想用于连续海流场：不确定集由多年误差数据标定，并用梯度直接搜索使选边翻转的流场。

\subsection{本文与已有工作的区别}
\label{sec:diff}
综上，已有工作分别解决了求最快航线、刻画海流的不确定性和获取新信息的问题，但没有回答一张海流图何时对某个具体决定不可信、何时值得付出观测的代价。表\ref{tab:innovation}逐项列出本文三个创新点与已有做法的区别、支持它们的实验证据，以及各自的适用条件。

\begin{table}[!htb]
\centering
\caption{原创创新点拆解：已有做法、本文做法、实验证据与适用条件。}
\label{tab:innovation}
\footnotesize
\setlength{\tabcolsep}{4pt}
\begin{tabular}{@{}>{\raggedright\arraybackslash}p{0.13\linewidth}>{\raggedright\arraybackslash}p{0.20\linewidth}>{\raggedright\arraybackslash}p{0.21\linewidth}>{\raggedright\arraybackslash}p{0.24\linewidth}>{\raggedright\arraybackslash}p{0.14\linewidth}@{}}
\toprule
创新点 & 已有做法及其不足 & 本文做法 & 实验证据 & 适用条件 \\
\midrule
1. 以合理反例判定该计算还是该观测 & 加密网格与元推理只减小求解误差，默认海流图准确；集合规划不回答何时转而观测。 & 以合理集合中的最大后悔值判定风险的来源：求解误差交给计算，粗图误差交给观测。 & 图\ref{fig:compute}：四种设置中计算量增加64倍，选错率的变化不超过2个百分点，换用当天0.5°图后降至7.0\%--10.4\%；81.0\%--99.7\%的决定存在合理反例。 & 单个禁航区的左右选边；单次任务中海流视为定常。 \\
\addlinespace
2. 数据标定的合理集合与包络定理梯度搜索 & 鲁棒方法的不确定集由分析者给定；集合方法需上千个成员才能估准，改选比例不计代价。 & 用多年误差标定误差大小、相关长度与半径；沿最快航线累加导数，每次迭代求解2次。 & 图\ref{fig:rank}、\ref{fig:budget}：约71次求解达到2048次求解的选边效果；前20\%的任务占53\%--71\%的额外航时；改用相似日差值构造集合后，保守改选不再有效（表\ref{tab:scope}）。 & 依赖开发年份的误差统计；流场变化大时一阶近似有误差。 \\
\addlinespace
3. 以最大后悔值决定何时观测 & 自适应采样为降低全场误差而测量；或一律观测，或从不观测。 & 最大后悔值超过阈值才观测，优先上浮取图；不能观测时改选最大后悔值较小的一侧。 & 图\ref{fig:respond}：湾流与阿古拉斯回流区的平均额外航时降至1.31与2.20小时；黑潮中保守改选从0.98降至0.70小时。 & 保守改选要求两侧最坏情形差别分明；黑潮中上浮代价需低于约0.8小时。 \\
\bottomrule
\end{tabular}
\end{table}
